%%Flux_bulk_simple
%%\[
 Les flux $H_{s}$, $L_{e}$ et $Tau$ sont calcul\'es pour chaque pas de temps.
%%\]
%Algorithme utilis� :
%
%En entr�e :
%Ps        hPa               Pressure on surface
%Uz        %                 Relative Humidity
%Ts        Kelvin or Celsius Air Temperature at height 0
%Tz        Kelvin or Celsius Air Temperature at height z
%Vz        m/s 	       Wind Strength (speed) at height z ; relatif a la mer si courant!!
%Ustar     m/s               Friction velocity 
%Hz        m                 Height of (Tz, Uz) measurements
%En sortie
%Hs        W/m2              Sensible heat flux 
%Le        W/m2              latent heat flux 
%Tau       N/m2              Turbulent part of wind stress 
%


\begin{align}
H_{s} & = -\rho_h  \times  C_{h}  \times  V_{z} \times C_{ph}  \times  (T_{sK} - T_{zK})\\
L_{e} & = -\rho_h  \times  C_{h}  \times  V_{z}  \times  L_{v}  \times  \frac{(r_{surf}\_sat - r_{z})}{1000}\\
Tau & =  \rho_h  \times  C_{d}  \times  {V_{z}}^{2}                                                
\end{align}

avec \\
D�finition des termes utilis�s :\\
$H_{z}$ : Hauteur de mesure en m (niveau z)\\
$R_{a}$ : Constante sp�cifique de l'air sec, en (J/K/kg) (=R/M$_{a}$=287.05)\\
$\rho_h$ : Masse volumique de l'air humide  ($kg/m^3$)\\
$C_{ph}$ : Chaleur massique (specific heat) � P constant de l'air humide (J/kg/K)\\
$C_{pa}$ : Chaleur massique (specific heat) � P constant de l'air sec (J/kg/K)\\
$T_{sK}$ : Temp�rature au niveau du sol (ou sst si c'est la mer) en Kelvin \\
$T_{zK}$ : Temp�rature au niveau z (z en m) en Kelvin \\
$T_{xK}$ signifie que la temp�rature est convertie en Kelvin si elle est fournie en celsius\\
$T_{v}$ : Temp�rature virtuelle en Kelvin au niveau z\\
$r_{surf}$ : rapport de m�lange en g/kg a la surface (sol ou mer)\\
$r_{surf}\_sat$ : rapport de m�lange satur� en g/kg a la surface (sol ou mer)\\
$r_{z}$ : rapport de m�lange en g/kg au niveau de mesure z\\
Le rapport de m�lange $r_{z}$ est calcul� par la fonction \texttt{rmel\_Hu()} \\
$V_{z}$ : Vitesse du vent en m/s au niveau z\\
$f_{m}$ : Fonction de stabilit� pour le mouvement d�pendant du signe du nombre de Richardson\\
$f_{h}$ : Fonction de stabilit� pour la chaleur et l'humidite d�pendant aussi du signe du nombre de Richardson\\
$C_{h}$ : Coefficient d'echange en surface pour la chaleur et l'humidite\\
$C_{d}$ : Coefficient d'echange en surface pour le mouvement\\
$cdn$ : Coefficient d'echange de quantit� de mouvement � la neutralit�\\
$K$ : Contante de Karman (entre 0.35 et 0.41)\\
$z0$ : Longueur de rugosit� en m ($\overline{u}$ est nul du sol jusqu'au niveau z0)\\
$z0$ est calcul�e ici par la formule de Charnock � partir de $ustar$ qui est la vitesse de vitesse de frottement en m/s.\\
$ustar$ est un des param�tres d'entr�e.\\
$Charnock$ : Coefficient valant 0.021\\
$g$ : Gravit� valant ici 9.81 $m/s^2$ (varie en fonction de la latitude)\\
$Ri$ : Nombre de Richardson calcul� ici avec la temp�rature virtuelle :
la vraie formule avec la temp�rature potentielle $\theta$ est :
\begin{eqnarray*}
 Ri & = & \dfrac{\frac{g}{\overline{\theta}} \times \frac{\overline{\theta_{z} - \theta_{0}}}{H_{z}}}{{(\frac{\overline{V_{z}}}{H_{z}})}^{2}} =
 \dfrac{g \times H_{z} \times (\overline{\theta_{z} - \theta_{0}})}{\overline{\theta}  \times  {\overline{V_{z}}^{2}}}
\end{eqnarray*}


\begin{eqnarray*}
\rho_h & = & P_{s}  \times  \frac{100}{R_{a}  \times  T_{v}}\\
C_{ph} & = & (C_{pa} + \frac{(\frac{r_{z}}{1000}  \times  Cpv)}{1 + (\frac{r_{z}}{1000})}\\
L_{v} & = & (2.501 - 0.00237  \times  (T_{sK} - 273.15) ) \times 10^6\\
z0 & = & Charnock \times \frac{{ustar}^{2}}{g}\\
cdn & = & \left( \frac{K}{log\frac{H_{z}}{z0}} \right) ^{2}\\
\text{on pose } dt  =  T_{v}-T_{sK}\\
\text{on calcule le nombre de Richardson} \\
Ri & = & g  \times  Hz  \times  \frac{dt}{(T_{v}  \times  {V_{z}}^{2})}\\
C_{h} & = & cdn  \times  f_{h}\\
C_{d} & = & cdn  \times  f_{m}\\
\end{eqnarray*}

%\usepackage{amsmath}
\begin{equation*}
\begin{aligned}
fm = &
\begin{cases}
1.-10. \times \dfrac{Ri}{(1.+75. \times cdn \times sqrt(\frac{H_{z}}{z0 \times fabs(Ri)}))} & \text{si $Ri <= 0$,}\\
\dfrac{1.}{(1.+10. \times \dfrac{Ri}{sqrt(1.+5. \times Ri))}} & \text{si $Ri > 0$.}\\
\end{cases} \\
& \text{avec } cdn=\frac{K}{{log(\frac{Hz}{z0})}^{2}}
\end{aligned}
\end{equation*}
\begin{equation*}
\begin{aligned}
fh = &
\begin{cases}
1.-15. \times \dfrac{Ri}{(1.+75. \times cdn \times sqrt(\frac{Hz}{z0 \times fabs(Ri)}))} & \text{si $Ri <= 0$,} \\
\dfrac{1.}{(1.+15. \times Ri \times sqrt(1.+5. \times Ri))} & \text{si $Ri > 0$.}\\
\end{cases} \\
\end{aligned}
\end{equation*}





